\section{Unit Tangent and Normal Vectors}\label{sec:tan_norm}

\subsection{Unit Tangent Vector}

Given a smooth vector-valued function \vrt, we defined in \autoref{def:vector_tangent} that any vector parallel to $\vrp(t_0)$ is \emph{tangent} to the graph of \vrt\ at $t=t_0$. It is often useful to consider just the \emph{direction} of $\vrp(t)$ and not its magnitude. Therefore we are interested in the unit vector in the direction of $\vrp(t)$. This leads to a definition.

\begin{definition}[Unit Tangent Vector]\label{def:unit_tangent}%
Let \vrt\ be a smooth function on an open interval $I$. The unit tangent vector $\vec T(t)$ is
\index{unit tangent vector!definition} \index{unit vector!unit tangent vector}
\[\vec T(t) = \frac1{\norm{\vrp(t)}}\vrp(t).\]
\end{definition}

\youtubeVideo{39LA5WyVgKY}{Tangent Line to a Parametrized Curve} % sic

\begin{example}[Computing the unit tangent vector]\label{ex_tannorm1}%
Let $\vrt =\bracket{3\cos t, 3\sin t, 4t}$. Find $\vec T(t)$ and compute $\vec T(0)$ and $\vec T(1)$.
\solution
We apply \autoref{def:unit_tangent} to find $\vec T(t)$. 
\begin{align*}
\vec T(t) &= \frac1{\norm{\vrp(t)}}\vrp(t) \\
				&=\frac{1}{\sqrt{\bigl(-3\sin t\bigr)^2+\bigl(3\cos t\bigr)^2+ 4^2}}\bracket{-3\sin t,3\cos t, 4}\\
				&=\bracket{-\frac35\sin t,\frac35\cos t,\frac45}.
\end{align*}
We can now easily compute $\vec T(0)$ and $\vec T(1)$:
%
\mtable{Plotting unit tangent vectors in \autoref{ex_tannorm1}.}{fig:tannorm1}{%
\myincludeasythree{
3Droll=0.4623610114058449,
3Dortho=0.004905222449451685,
3Dc2c=0.6268976330757141 0.5422312617301941 0.5594503283500671,
3Dcoo=-3.587019205093384 4.6990532875061035 15.216367721557617,
3Droo=150.00000240039893}{\marginparwidth}{Blue space‐curve 𝑟(𝑡)=⟨t, t^3, 10 sin t⟩ red unit-tangent 𝑇 shown at two marked t values.}{figures/figtannorm1_3D}}
%
\[
\vec T(0) =\bracket{0,\frac35,\frac45}\,;\quad \vec T(1) =\bracket{-\frac35\sin 1,\frac35\cos 1,\frac45}%\approx\bracket{-0.505,0.324,0.8}
.
\]
These are plotted in \autoref{fig:tannorm1} with their initial points at $\vec r(0)$ and $\vec r(1)$, respectively. (They look rather ``short'' since they are only length 1.)
%
%The unit tangent vector $\vec T(t)$ always has a magnitude of 1, though it is sometimes easy to doubt that is true. We can help solidify this thought in our minds by computing $\norm{\vec T(1)}$:
%\[\norm{\vec T(1)} \approx \sqrt{(-0.505)^2+0.324^2+0.8^2} = 1.000001.\]
%We have rounded in our computation of $\vec T(1)$, so we don't get 1 exactly. We leave it to the reader to use the exact representation of $\vec T(1)$ to verify it has length 1.
\end{example}

In many ways, the previous example was ``too nice.'' It turned out that $\vrp(t)$ was always of length 5. In the next example the length of $\vrp(t)$ is variable, leaving us with a formula that is not as clean.

\begin{example}[Computing the unit tangent vector]\label{ex_tannorm2}%
Let $\vrt=\bracket{t^2-t,t^2+t}$. Find $\vec T(t)$ and compute $\vec T(0)$ and $\vec T(1)$.
\solution
We find $\vrp(t) =\bracket{2t-1,2t+1}$, and
%
\mtable{Plotting unit tangent vectors in \autoref{ex_tannorm2}.}{fig:tannorm2}{\begin{tikzpicture}[alt={Plane curve r(t)=⟨t², t³, 3t⟩ three red unit-tangent T(t), drawn where t = -1, 0, 2.},>=stealth]
\begin{axis}[width=1.16\marginparwidth,tick label style={font=\scriptsize},
axis y line=middle,axis x line=middle,name=myplot,
ymin=-1.1,ymax=6.5,xmin=-2.1,xmax=7]
\addplot [thick,draw={\colorone}, smooth,domain=-2:2,samples=30] ({x^2-x},{x^2+x});
\draw [thick,->,draw={\colortwo}] (axis cs:0,0)--(axis cs:-.707,.707);
\draw [thick,->,draw={\colortwo}] (axis cs:0,2)--(axis cs:0.32,2.95);
\draw [thick,->,draw={\colorone}] (axis cs:3.79,.77)--(axis cs:3.75,.75);
\end{axis}
\node [right] at (myplot.right of origin) {\scriptsize $x$};
\node [above] at (myplot.above origin) {\scriptsize $y$};
\end{tikzpicture}}
%
\[\norm{\vrp(t)} = \sqrt{(2t-1)^2+(2t+1)^2} = \sqrt{8t^2+2}.\]
Therefore
\[\vecT(t) = \frac{1}{\sqrt{8t^2+2}}\bracket{2t-1,2t+1}=\bracket{\frac{2t-1}{\sqrt{8t^2+2}},\frac{2t+1}{\sqrt{8t^2+2}}}.\]
When $t=0$, we have $\vec T(0) =\bracket{-1/\sqrt{2},1/\sqrt{2}}$; when $t=1$, we have $\vec T(1) =\bracket{1/\sqrt{10}, 3/\sqrt{10}}.$ We leave it to the reader to verify each of these is a unit vector. They are plotted in \autoref{fig:tannorm2}.
\end{example}

\subsection{Unit Normal Vector}

\mtable{Given a direction in the plane, there are always two directions orthogonal to it.}{fig:tannorm_intro1}{\begin{tikzpicture}[alt={At a single point on the previous curve, red T
and black unit normals t_1, t_2, illustrate the two orthogonal choices.},>=stealth]
\begin{axis}[width=1.16\marginparwidth,tick label style={font=\scriptsize},
axis y line=middle,axis x line=middle,name=myplot,xtick=\empty,ytick=\empty,
ymin=-.1,ymax=1.9,xmin=-.1,xmax=2.3]
\addplot [thick,draw={\colorone}, smooth,domain=-1:1,samples=30] ({x^2-x+.6},{x^2+x+.3});
\draw[thick,->,black] (axis cs:.51,.41) -- (axis cs:.233,.83);
\draw[thick,->,draw={\colortwo}] (axis cs:.51,.41) -- (axis cs:.1,0.13);
\draw[thick,->,draw={\colortwo}] (axis cs:.51,.41) -- (axis cs:.93,.69);
\draw [thick,->,draw={\colorone}] (axis cs:0.71, 0.21)--(axis cs:.698,.218);
\end{axis}
\node [right] at (myplot.right of origin) {\scriptsize $x$};
\node [above] at (myplot.above origin) {\scriptsize $y$};
\end{tikzpicture}}

Just as knowing the direction tangent to a path is important, knowing a direction orthogonal to a path is important. When dealing with real-valued functions, we defined the normal line at a point to the be the line through the point that was perpendicular to the tangent line at that point. We can do a similar thing with vector-valued functions. Given \vrt\ in $\mathbb{R}^2$, we have 2 directions perpendicular to the tangent vector, as shown in \autoref{fig:tannorm_intro1}. It is good to wonder ``Is one of these two directions preferable over the other?''

Given \vrt\ in $\mathbb{R}^3$,  there are infinite vectors orthogonal to the tangent vector at a given point. Again, we might wonder ``Is one of these infinite choices preferable over the others? Is one of these the `right' choice?''

The answer in both $\mathbb{R}^2$ and $\mathbb{R}^3$ is ``Yes, there is one vector that is preferable.'' %, it is the `right' one to choose.''
Recall \autoref{thm:vects_of_constant_length}, which states that if \vrt\ has constant length, then \vrt\ is orthogonal to $\vrp(t)$ for all $t$. We know $\vec T(t)$, the unit tangent vector, has constant length. Therefore $\vec T(t)$ is orthogonal to $\vec T\,'(t)$.\index{vector-valued function!of constant length}

We'll see that $\vec T\,'(t)$ is more than just a convenient choice of vector that is orthogonal to $\vrp(t)$; rather, it is the ``right'' choice. Since all we care about is the direction, we define this newly found vector to be a unit vector.

\mnote{\textbf{Note:} $\vec T(t)$ is a unit vector, by definition. This \emph{does not} imply that $\vec T\,'(t)$ is also a unit vector.}

\begin{definition}[Unit Normal Vector]\label{def:unit_normal}%
Let \vrt\ be a vector-valued function where the unit tangent vector, $\vec T(t)$, is smooth on an open interval $I$. The \textbf{unit normal vector} $\vec N(t)$ is
\index{unit normal vector!definition}\index{unit vector!unit normal vector}
\[\vec N(t) = \frac1{\norm{\vec T\,'(t)}}\vec T\,'(t).\]
\end{definition}

\begin{example}[Computing the unit normal vector]\label{ex_tannorm3}%
Let $\vrt =\bracket{3\cos t, 3\sin t, 4t}$ as in \autoref{ex_tannorm1}. Sketch both $\vec T(\pi/2)$ and $\vec N(\pi/2)$ with initial points at $\vec r(\pi/2)$.
\solution
In \autoref{ex_tannorm1}, we found $\vec T(t)=\bracket{-\frac35\sin t,\frac35\cos t,4/5}$. Therefore 
\[\vec T\,'(t) =\bracket{-\frac35\cos t,-\frac35\sin t,0}\quad \text{and} \quad \norm{\vec T\,'(t)} = \frac35.\]
%
\mtable{Plotting unit tangent and normal vectors in \autoref{ex_tannorm3}.}{fig:tannorm3}{%
\myincludeasythree{
3Droll=0.9552187107340909,
3Dortho=0.005093402229249477,
3Dc2c=0.5105816125869751 0.5339998006820679 0.6739070415496826,
3Dcoo=4.143548488616943 1.2702205181121826 14.294137954711914,
3Droo=150.0000014005964}{\marginparwidth}{Along the helix, red T and black N form the Frenet pair at t = 1.2; axes ticked at ±3, ±10.}{figures/figtannorm3_3D}}%
%
Thus
\[\vec N(t) = \frac{\vec T\,'(t)}{3/5} =\bracket{-\cos t,-\sin t,0}.\]
We compute $\vec T(\pi/2) =\bracket{-3/5,0,4/5}$ and $\vec N(\pi/2) =\bracket{0,-1,0}$. These are sketched in \autoref{fig:tannorm3}.
\end{example}

The previous example was once again ``too nice.'' In general, the expression for $\vec T(t)$ contains fractions of square-roots, hence the expression of $\vec T\,'(t)$ is very messy. We demonstrate this in the next example.

\begin{example}[Computing the unit normal vector]\label{ex_tannorm4}%
Let $\vrt=\bracket{t^2-t,t^2+t}$ as in \autoref{ex_tannorm2}. Find $\vec N(t)$ and sketch $\vrt$ with the unit tangent and normal vectors at $t=-1,0$ and 1.
\solution
In \autoref{ex_tannorm2}, we found
\[\vec T(t) =\bracket{\frac{2t-1}{\sqrt{8t^2+2}},\frac{2t+1}{\sqrt{8t^2+2}}}.\]
Finding $\vec T\,'(t)$ requires two applications of the Quotient Rule:
%
\begin{align*}
\vec T\,'(t) &=
	\bracket{\frac{\sqrt{8t^2+2}(2)-(2t-1)\left(\frac12(8t^2+2)^{-1/2}(16t)\right)}{8t^2+2},\right.\\
	&\qquad\qquad\left.\frac{\sqrt{8t^2+2}(2)-(2t+1)\left(\frac12(8t^2+2)^{-1/2}(16t)\right)}{8t^2+2}}\\
	&=\bracket{\frac{4 (2 t+1)}{\left(8 t^2+2\right)^{3/2}},
	\frac{4(1-2 t)}{\left(8 t^2+2\right)^{3/2}}}
\end{align*}

This is not a unit vector; to find $\vec N(t)$, we need to divide $\vec T\,'(t)$ by it's magnitude.
\begin{align*}
\norm{\vec T\,'(t)} &= \sqrt{\frac{16(2t+1)^2}{(8t^2+2)^3}+\frac{16(1-2t)^2}{(8t^2+2)^3}}\\
					&= \sqrt{\frac{16(8t^2+2)}{(8t^2+2)^3}}\\
					&= \frac{4}{8t^2+2}.
\end{align*}
Finally, 
\begin{align*}
\vec N(t) &= \frac1{4/(8t^2+2)}\bracket{\frac{4 (2 t+1)}{\left(8 t^2+2\right)^{3/2}},\frac{4
   (1-2 t)}{\left(8 t^2+2\right)^{3/2}}}\\
	&=\bracket{\frac{2t+1}{\sqrt{8t^2+2}},-\frac{2t-1}{\sqrt{8t^2+2}}}.
\end{align*}
Because we are normalizing $\vec T\,'(t)$, it is usually easier to scale it first.  We see that $\vec T\,'(t)$ is parallel to $\bracket{2t+1,1-2t}$, which has length $\sqrt{(2t+1)^2+(1-2t)^2}=\sqrt{8t^2+2}$, leading to the same $\vec N(t)$.

\mtable{Plotting unit tangent and normal vectors in \autoref{ex_tannorm4}.}{fig:tannorm4}{\begin{tikzpicture}[alt={Plane curve; red T and N shown at three equally spaced parameter values.},>=stealth]
\begin{axis}[width=1.16\marginparwidth,tick label style={font=\scriptsize},
axis y line=middle,axis x line=middle,name=myplot,
ymin=-1.1,ymax=6.5,xmin=-2.1,xmax=7]
\addplot [thick,draw={\colorone}, smooth,domain=-2:2,samples=30] ({x^2-x},{x^2+x});
\draw [thick,->,draw={\colortwo}] (axis cs:0,0)--(axis cs:-.707,.707);
\draw [thick,->,draw={\colortwo}] (axis cs:0,2)--(axis cs:0.32,2.95);
\draw [thick,->,draw={\colortwo}] (axis cs:2,0)--(axis cs:1.05,-.32);
\draw [thick,->,draw={\colortwo}] (axis cs:0,0)--(axis cs:.707,.707);
\draw [thick,->,draw={\colortwo}] (axis cs:0,2)--(axis cs:.95,1.68);
\draw [thick,->,draw={\colortwo}] (axis cs:2,0)--(axis cs:1.68,0.95);
\draw [thick,->,draw={\colorone}] (axis cs:3.79,.77)--(axis cs:3.75,.75);
\end{axis}
\node [right] at (myplot.right of origin) {\scriptsize $x$};
\node [above] at (myplot.above origin) {\scriptsize $y$};
\end{tikzpicture}}

Using this formula for $\vec N(t)$, we compute the unit tangent and normal vectors for $t=-1,0$ and 1 and sketch them in \autoref{fig:tannorm4}.
\end{example}

The final result for $\vec N(t)$ in \autoref{ex_tannorm4} is suspiciously similar to $\vec T(t)$. There is a clear reason for this. If $\vec u =\bracket{u_1,u_2}$ is a unit vector in $\mathbb{R}^2$, then the \emph{only} unit vectors orthogonal to $\vec u$ are $\bracket{-u_2,u_1}$ and $\bracket{u_2,-u_1}$. Given $\vec T(t)$, we can quickly determine $\vec N(t)$ if we know which term to multiply by $(-1)$.

Consider again \autoref{fig:tannorm4}, where we have plotted some unit tangent and normal vectors. Note how $\vec N(t)$ always points ``inside'' the curve, or to the concave side of the curve. This is not a coincidence; this is true in general. Knowing the direction that $\vec r(t)$ ``turns'' allows us to quickly find $\vec N(t)$.


\begin{theorem}[Unit Normal Vectors in $\mathbb{R}^2$]\label{thm:concave}%
Let $\vec r(t)$ be a vector-valued function in $\mathbb{R}^2$ where $\vec T\,'(t)$ is smooth on an open interval $I$. Let $t_0$ be in $I$ and $\vec T(t_0) =\bracket{t_1,t_2}$ Then $\vec N(t_0)$ is either
\[\vec N(t_0) =\bracket{-t_2,t_1}\quad \text{or}\quad \vec N(t_0) =\bracket{t_2,-t_1},\]
whichever is the vector that points to the concave side of the graph of $\vec r$.
\index{unit tangent vector!in $\mathbb{R}^2$}\index{unit normal vector!in $\mathbb{R}^2$}
\end{theorem}

\subsection{Application to Acceleration}

Let \vrt\ be a position function. It is a fact (stated later in \autoref{thm:acc_plane}) % and proved in the Exercises) 
 that acceleration, \vat, lies in the plane defined by $\vec T$ and $\vec N$. That is, there are scalars $a_{\mathrm{T}}$ and $a_{\mathrm{N}}$ such that 
\[\vat = a_{\mathrm{T}}\vec T(t) + a_{\mathrm{N}}\vec N(t).\]
The scalar $a_{\mathrm{T}}$ measures ``how much'' acceleration is in the direction of travel, that is, it measures the component of acceleration that affects the speed. The scalar $a_{\mathrm{N}}$ measures ``how much'' acceleration is perpendicular to the direction of travel, that is, it measures the component of acceleration that affects the direction of travel.
\index{unit tangent vector!and acceleration}\index{unit normal vector!and acceleration}

We can find $a_{\mathrm{T}}$ using the orthogonal projection of $\vec a(t)$ onto $\vec T(t)$ (review \autoref{def:orthogonal_projection} in \autoref{sec:dot_product} if needed).
Recalling that since $\vec T(t)$ is a unit vector, $\vec T(t)\cdot\vec T(t)=1$, so we have 
\[\proj{a(t)}{T(t)} = \frac{\vec a(t)\cdot\vec T(t)}{\vec T(t)\cdot\vec T(t)}\vec T(t) = \underbrace{\bigl(\vec a(t)\cdot\vec T(t)\bigr)}_{a_{\mathrm{T}}}\vec T(t).\]
Thus the amount of \vat\ in the direction of $\vec T(t)$ is $a_{\mathrm{T}}=\vat\cdot\vec T(t)$. The same logic gives $a_{\mathrm{N}} = \vat\cdot\vec N(t)$.

While this is a fine way of computing $a_{\mathrm{T}}$, there are simpler ways of finding $a_{\mathrm{N}}$ (as finding $\vec N$ itself can be complicated). The following theorem gives alternate formulas for $a_{\mathrm{T}}$ and $a_{\mathrm{N}}$.

\mnote{\textbf{Note:} Keep in mind that both $a_{\mathrm{T}}$ and $a_{\mathrm{N}}$ are functions of $t$; that is, the scalar changes depending on $t$. It is convention to drop the ``$(t)$'' notation from $a_{\mathrm{T}}(t)$ and simply write $a_{\mathrm{T}}$.}

{
\tcbset{grow to right by=2em}
\begin{theorem}[Acceleration in the Plane Defined by $\vec T$ and $\vec N$]\label{thm:acc_plane}%
Let \vrt\ be a position function with acceleration \vat\ and unit tangent and normal vectors $\vec T(t)$ and $\vec N(t)$. Then $\vat$ lies in the plane defined by $\vec T(t)$ and $\vec N(t)$; that is, there exists scalars $a_{\mathrm{T}}$ and $a_{\mathrm{N}}$ such that
\[\vat = a_{\mathrm{T}}\vec T(t) + a_{\mathrm{N}}\vec N(t).\]
Moreover,
\begin{align*}
a_\mathrm{T} &= \vat\cdot\vec T(t) = \frac{\dd}{\dd t}\Bigl(\norm{\vec v(t)}\Bigr) \\
a_\mathrm{N} &= \vat\cdot \vec N(t) = \sqrt{\norm{\vec a(t)}^2-a_\mathrm{T}^2} = \frac{\norm{\vat\times\vvt}}{\norm{\vvt}} = \norm{\vvt}\,\norm{\vec T\,'(t)}
\end{align*}%
\index{unit tangent vector!and acceleration}%
\index{unit normal vector!and acceleration}%
\index{an@$a_{\mathrm{N}}$}\index{unit normal vector!an@$a_{\mathrm{N}}$}%
\index{at@$a_{\mathrm{T}}$}\index{unit tangent vector!at@$a_{\mathrm{T}}$}%
\end{theorem}
}
% GregH asked about this at https://tex.stackexchange.com/q/176439/107497
% egreg provided a solution at https://tex.stackexchange.com/a/88357/107497
% but LaTeXML sends \string to the index, which messes up the output
% this appears to now be fixed

Note the second formula for $a_{\mathrm{T}}$: $\ds \frac{\dd}{\dd t}\Bigl(\norm{\vvt}\Bigr)$. This measures the rate of change of speed, which again is the amount of acceleration in the direction of travel.

\begin{proof}
We see that
\begin{align*}
 \vat
 &=\frac\dd{\dd t}\vvt=\frac\dd{\dd t}\left(\norm{\vvt}\,\vec T(t)\right)
 =\left(\frac\dd{\dd t}\norm{\vvt}\right)\,\vec T(t)+\norm{\vvt}\vec T\,'(t) \\
 &=\left(\frac\dd{\dd t}\norm{\vvt}\right)\,\vec T(t)+\norm{\vvt}\norm{\vec T\,'(t)}\vec N(t).
\end{align*}
Since $\vec T(t)$ and $\vec N(t)$ are not parallel, this decomposition is unique and the coefficients tell us $a_\mathrm{T}$ and $a_\mathrm{N}$.

Because $\norm{\vec T}=1$, \autoref{thm:vects_of_constant_length} tells us that $\vec T$ and $\vec T\,'=\norm{\vec T\,'}\vec N$ are orthogonal.  This means that
\[
 \norm{\vat\times\vvt}
 =\norm{a_\mathrm{N}\vec N(t)\times\norm{\vvt}\,\vec T(t)}
 =a_\mathrm{N}\norm{\vvt}.
\]
Also, the Pythagorean theorem tells us that
\[
 \norm{\vec a(t)}^2=\norm{a_\mathrm{T}\vec T(t)}^2+\norm{a_\mathrm{N}\vec N(t)}^2
 =a_\mathrm{T}^2+a_\mathrm{N}^2.\qedhere
\]
\end{proof}

\begin{example}[Computing $a_\mathrm{T}$ and $a_\mathrm{N}$]\label{ex_tannorm5}%
Let $\vrt =\bracket{3\cos t, 3\sin t, 4t}$ as in Examples \ref{ex_tannorm1} and \ref{ex_tannorm3}. Find $a_\mathrm{T}$ and $a_\mathrm{N}$.
\solution
The previous examples give $\vat =\bracket{-3\cos t,-3\sin t,0}$ and 
\[\vec T(t) =\bracket{-\frac35\sin t,\frac35\cos t,\frac45}\quad \text{and}\quad \vec N(t) =\bracket{-\cos t,-\sin t,0}.\]
We can find $a_\mathrm{T}$ and $a_\mathrm{N}$ directly with dot products:
\begin{align*}
a_\mathrm{T} &= \vat\cdot \vec T(t) = \frac95\cos t\sin t-\frac95\cos t\sin t+0 = 0.\\
a_\mathrm{N} &= \vat\cdot \vec N(t) = 3\cos^2t+3\sin^2t + 0 = 3.
\end{align*}
Thus $\vat = 0\vec T(t) + 3\vec N(t) = 3\vec N(t)$, which is clearly the case.

What is the practical interpretation of these numbers? $a_\mathrm{T}=0$ means the object is moving at a constant speed, and hence all acceleration comes in the form of direction change.
\end{example}

\begin{example}[Computing $a_\mathrm{T}$ and $a_\mathrm{N}$]\label{ex_tannorm6}%
Let $\vrt=\bracket{t^2-t,t^2+t}$ as in Examples \ref{ex_tannorm2} and \ref{ex_tannorm4}. Find $a_\mathrm{T}$ and $a_\mathrm{N}$.
\solution
The previous examples give $\vat =\bracket{2,2}$ and
\[\vec T(t) =\bracket{\frac{2t-1}{\sqrt{8t^2+2}},\frac{2t+1}{\sqrt{8t^2+2}}}\quad \text{and}\quad \vec N(t) =\bracket{\frac{2t+1}{\sqrt{8t^2+2}},-\frac{2t-1}{\sqrt{8t^2+2}}}.\]
While we can compute $a_\mathrm{N}$ using $\vec N(t)$, we instead demonstrate using another formula from \autoref{thm:acc_plane}.
\begin{align*}
a_\mathrm{T} &= \vat\cdot\vec T(t) = \frac{4t-2}{\sqrt{8t^2+2}}+\frac{4t+2}{\sqrt{8t^2+2}} = \frac{8t}{\sqrt{8t^2+2}}.\\
a_\mathrm{N} &= \sqrt{\norm{\vat}^2-a_\mathrm{T}^2} = \sqrt{8-\left(\frac{8t}{\sqrt{8t^2+2}}\right)^2}=\frac{4}{\sqrt{8t^2+2}}.\medskip
\end{align*}
When $t=2$, $\ds a_{\mathrm{T}} = \frac{16}{\sqrt{34}}%\approx 2.74
$ and $\ds a_{\mathrm{N}} = \frac{4}{\sqrt{34}}% \approx 0.69
$. We interpret this to mean that at $t=2$, the particle is acculturating mostly by increasing speed, not by changing direction. As the path near $t=2$ is relatively straight, this should make intuitive sense. \autoref{fig:tannorm6} gives a graph of the path for reference.
\mtable{Graphing $\vec r(t)$ in \autoref{ex_tannorm6}.}{fig:tannorm6}{\begin{tikzpicture}[alt={Curve r(t)=(t2,t3−3t⟩r(t)=⟨t2 ,t3 −3t⟩; orientation arrows added, endpoints labelled t=0 and t
= 2.},>=stealth]
\begin{axis}[width=1.16\marginparwidth,tick label style={font=\scriptsize},
axis y line=middle,axis x line=middle,name=myplot,
ymin=-1.1,ymax=6.5,xmin=-2.1,xmax=7]
\addplot [thick,draw={\colorone}, smooth,domain=-2.1:2.1,samples=30] ({x^2-x},{x^2+x});
\filldraw [black] (axis cs: 2,6) circle (1.5pt)
 node [below right] {\scriptsize $t=2$};
\filldraw [black] (axis cs: 0,0) circle (1.5pt)
 node [above right] {\scriptsize $t=0$};
\draw (axis cs: 4.2,1.8) node {\scriptsize $\vec r(t)$};
\draw [thick,->,draw={\colorone}] (axis cs:3.79,.77)--(axis cs:3.75,.75);
\end{axis}
\node [right] at (myplot.right of origin) {\scriptsize $x$};
\node [above] at (myplot.above origin) {\scriptsize $y$};
\end{tikzpicture}}

Contrast this with $t=0$, where $a_{\mathrm{T}}=0$ and $a_{\mathrm{N}}=4/\sqrt2=2\sqrt2 %\approx 2.82
$. Here the particle's speed is not changing and all acceleration is in the form of direction change.
\end{example}

\begin{example}[Analyzing projectile motion]\label{ex_tannorm7}%
A ball is thrown from a height of 240ft with an initial speed of 64ft/s and an angle of elevation of $30^\circ$. Find the position function \vrt\ of the ball and analyze $a_\mathrm{T}$ and $a_\mathrm{N}$.\index{projectile motion}
\solution
Using \autoeqref{idea:projectile} of \autoref{sec:vvf_motion} we form the position function of the ball:
\[
\vrt =\bracket{\bigl(64\cos 30^\circ\bigr)t, -16t^2+\bigl(64\sin 30^\circ\bigr)t+240},
\]
which we plot in \autoref{fig:tannorm7b}.

\mtable{Plotting the position of a thrown ball, with 1s increments shown.}{fig:tannorm7b}{\begin{tikzpicture}[alt={Projectile arc: 1s dots from launch to t=5; feet scale, max 𝑥=300, peak 
y=220.},>=stealth]
\begin{axis}[width=\marginparwidth,tick label style={font=\scriptsize},
axis y line=middle,axis x line=middle,name=myplot,
ymin=-10,ymax=260,xmin=-10,xmax=310]
\addplot [thick,draw={\colorone}, smooth,domain=0:5,samples=30]
 ({64*cos(30)*x},{-16*x^2+64*sin(30)*x+240});
\filldraw [black] (axis cs: 0,240) circle (1.5pt)
 node [below right] {\scriptsize $t=0$};
\filldraw [black] (axis cs: 55.4,256) circle (1.5pt)
 node (A)[ ] {};
\filldraw [black] (axis cs: 111,240) circle (1.5pt)
 node [above right] {\scriptsize $t=2$};
\filldraw [black] (axis cs: 166,192) circle (1.5pt)
 node [above right] {\scriptsize $t=3$};
\filldraw [black] (axis cs: 222,112) circle (1.5pt)
 node [above right] {\scriptsize $t=4$};
\filldraw [black] (axis cs: 277,0) circle (1.5pt) node (B) [] {};
\draw [thick,->,draw={\colorone}] (axis cs:149.649, 209.76)--(axis cs:152,207);
\end{axis}
\node [right] at (myplot.right of origin) {\scriptsize $x$};
\node [above] at (myplot.above origin) {\scriptsize $y$};
\draw (A) node[above] {{\scriptsize $t=1$}};
\draw (B) node [shift={(9pt,8pt)}] {\scriptsize $t=5$};
\end{tikzpicture}}

From this we find
\[
 \vvt =\bracket{64\cos 30^\circ, -32t+64\sin 30^\circ}
 \qquad and\qquad
 \vat =\bracket{0,-32}.
\]
Computing $\vec T(t)$ is not difficult, and with some simplification we find
\[
 \vec T(t) =\bracket{\frac{\sqrt{3}}{\sqrt{t^2-2t+4}}, \frac{1-t}{\sqrt{t^2-2t+4}}}.
\]
With $\vat$ as simple as it is, finding $a_\mathrm{T}$ is also  simple:
\[a_\mathrm{T} = \vat\cdot \vec T(t) = \frac{32t-32}{\sqrt{t^2-2t+4}}.\]
We skip finding $\vec N(t)$ and find $a_\mathrm{N}$ through the formula $a_\mathrm{N} = \sqrt{\norm{\vat}^2-a_\mathrm{T}^2\,}$\,:
\[a_\mathrm{N} = \sqrt{32^2-\left(\frac{32t-32}{\sqrt{t^2-2t+4}}\right)^2} = \frac{32\sqrt{3}}{\sqrt{t^2-2t+4}}.\]
\autoref{fig:tannorm7a} gives a table of values of $a_\mathrm{T}$ and $a_\mathrm{N}$. When $t=0$, we see the ball's speed is decreasing; when $t=1$ the speed of the ball is unchanged. This corresponds to the fact that at $t=1$ the ball reaches its highest point.

After $t=1$ we see that $a_\mathrm{N}$ is decreasing in value. This is because as the ball falls, it's path becomes straighter and most of the acceleration is in the form of speeding up the ball, and not in changing its direction.
\mtable{A table of values of $a_\mathrm{T}$ and $a_\mathrm{N}$ in \autoref{ex_tannorm7}.}{fig:tannorm7a}{%
\tagpdfsetup{table/header-rows={1}}
\begin{tabular}{cll}
$t$ & $a_\mathrm{T}$ & $a_\mathrm{N}$ \\ \midrule
0 & $-16$ & 27.7 \\
1 & $\phantom{-0}0$ & 32 \\
2 & $\phantom{-}16$ & 27.7\\
3 & $\phantom{-}24.2$ & 20.9 \\
4 & $\phantom{-}27.7$ & 16 \\
5 & $\phantom{-}29.4$ & 12.7 
\end{tabular}}
\end{example}

Our understanding of the unit tangent and normal vectors is aiding our understanding of motion. The work in \autoref{ex_tannorm7} gave quantitative analysis of what we intuitively knew.

The next section provides two more important steps towards this analysis. %One problem we currently have is that while we can accurately describe a path for a particle, we are not particularly good at regulating the speed at which the particle travels. For instance, consider the parabolic path of \autoref{ex_tannorm2}. Our particle slows as it nears the vertex, but clearly it ``doesn't have to.'' What vector-valued function has the same path, but where the particle speeds up near the vertex? 
We currently describe position only in terms of time. In everyday life, though, we often describe position in terms of distance (``The gas station is about 2 miles ahead, on the left.''). The \emph{arc length parameter} allows us to reference position in terms of distance traveled. 

We also intuitively know that some paths are straighter than others --- and some are curvier than others, but we lack a measurement of ``curviness.'' The arc length parameter provides a way for us to compute \emph{curvature}, a quantitative measurement of how curvy a curve is.

\printexercises{exercises/11-04-exercises}
